\scanpage{6}
base sur $W$ formée des
\[
\Theta_0^{a_0}\cdots\Theta_\nu^{a_\nu}\qquad 0\leq a_i\leq p-1\quad\text{pour }i=0,\ldots,\nu.
\]
Il faut maintenant diviser par l'idéal \inserted{$I$} engendré par $\Theta_0$. On a le

\uline{Lemme} Pour que $\displaystyle\sum_{\substack{a=(a_0,\ldots a_\nu)\\0\leq a_i\leq p-1}}w_a.\Theta_0^{a_0}\cdots\Theta_\nu^{a_\nu}$ soit dans l'idéal engendré par $\Theta_0$, il faut et il suffit que

a) pour $\forall\ a=(a_0,\ldots a_\nu)$, si \inserted{$a\ne0$ et si} $a_i$ est la première composante non nulle de $a$, on ait $\deleted{\unsure{p}}w_a\in p^iW$

b) le coefficient $w_0$ de $1$ soit dans $p^{c(\nu)}W$, où
\[
c(\nu)=\nu+p[p^\nu-(p^{\nu-1}+p^{\nu-2}+\cdots+1)]
\]

\uline{Dém.} \uline{Condition suffisante} Résulte des relations
\[
\left\{\begin{aligned}p^i\Theta_i&\in I\\p^{c(\nu)}&\in I\end{aligned}\right.
\]
La première relation par récurrence sur $i\geq0$. Si $i=0$ c'est $\Theta_0\in I$, \deleted{\unsure{c'est le}} trivial. Supposons qu'on sache $p^i\Theta_i\in I$ $(i\leq\nu-1)$ on utilise
\[
(\Theta_i+\lambda_i)^p=\frac{p^{i+1}!}{(p^i!)^p}(\Theta_{i+1}+\lambda_{i+1}),
\]
qui s'écrit
\[
\deleted{\unsure{\Theta_i^p=}}\quad\sum_{\alpha=1}^{\unsure{p-1|p}}\inserted{C_i^p}\Theta_i^\alpha\lambda_i^{p-\alpha}
=\left(\deleted{-\lambda_i^p}+\deleted{\frac{p^{i+1}!}{(p^i!)^p}\unsure{\lambda_{i+1}|\Theta_{i+1}}}+\frac{p^{i+1}!}{(p^i!)^p}\Theta_{i+1}\right)
\]
multipliant par $p^i$, on trouve que, compte tenu que $p^i\Theta_i\in I$ donc $p_i\Theta_i^\alpha=0$ si $\alpha\geq1$ :
\[
p^i\frac{p^{i+1}!}{(p^i!)^p}\Theta_{i+1}=0
\]
Or on constate que
\[
v_p\left(\frac{p^{i+1}!}{(p^i!)^p}\right)=1
\]
\mspage{5}
